\section{The full zeta crosses its first circle}
\label{sec:global}

The natural-boundary sector in Theorem~\ref{thm:witness} does not determine
the first boundary of the full product.  Define the total local correction
\[
 \Delta_n=\Ninf_n-\Nsol_n.
 \tag{6.1}\label{6.1}
\]
By Theorem~\ref{thm:parity}, only the \(L_n\) cyclic no-\texttt{aa} words
contribute.  For such a word, \(D_w\) is even and
\[
 \frac{D_w}{2}
 \le D_w-\oddpart(D_w)<D_w.
 \tag{6.2}\label{6.2}
\]

\begin{theorem}[Exact correction scale]
\label{thm:correction}
Let \(\golden=(1+\sqrt5)/2\).  Then
\[
 \frac12(8^n-2\beta^n+1)L_n
 \le\Delta_n\le(8^n+1)L_n
 \tag{6.3}\label{6.3}
\]
for all sufficiently large \(n\), and consequently
\[
 \lim_{n\to\infty}\Delta_n^{1/n}=8\golden.
 \tag{6.4}\label{6.4}
\]
\end{theorem}

\begin{proof}
Sum \eqref{6.2} over the active words and apply the uniform determinant
bounds \eqref{3.5}.  Since \(\beta<8\) and \(L_n^{1/n}\to\golden\), taking
\(n\)-th roots in \eqref{6.3} proves \eqref{6.4}.
\end{proof}

Set
\[
 G(z)=\sum_{n\ge1}\frac{\Delta_n}{n}z^n,
 \qquad R=\frac1{8\golden}\approx0.077254.
 \tag{6.5}\label{6.5}
\]
Theorem~\ref{thm:correction} makes \(R\) the exact radius of \(G\).  In the
initial disk of convergence,
\[
 \Zsol(z)=\Zinf(z)e^{-G(z)}.
 \tag{6.6}\label{6.6}
\]

\begin{theorem}[Pole-removed zero-free disk]
\label{thm:continuation}
The function \(\Zsol\) has a unique meromorphic continuation to \(|z|<R\).
In this disk it has exactly one simple pole, at \(z=1/16\), and no zeros.
Equivalently,
\[
 H(z):=(1-16z)\Zsol(z)
 =\frac{1-13z+33z^2}{1-2z}e^{-G(z)}
 \tag{6.7}\label{6.7}
\]
is holomorphic and nowhere zero on \(|z|<R\).  The Taylor series of
\(\Zsol\) at the origin has radius exactly \(1/16\), but its first
convergence circle is not a natural boundary.
\end{theorem}

\begin{proof}
The series \(G\) is holomorphic in \(|z|<R\), so \(e^{-G}\) is holomorphic
and zero-free there.  Equation \eqref{6.6} therefore defines the unique
meromorphic continuation from the initial disk.  The numerator zeros of
\(\Zinf\) have moduli
\[
 \frac{2}{13+\sqrt{37}}\approx0.104806,
 \qquad
 \frac{2}{13-\sqrt{37}}\approx0.289133,
\]
both greater than \(R\).  Of the denominator zeros, only \(1/16\) lies in
the disk, and it is simple.  Nothing cancels it.  Formula \eqref{6.7} and all
assertions follow.
\end{proof}

Positive coefficients and the lower bound in \eqref{6.3} imply
\(G(r)\to+\infty\) as \(r\uparrow R\).  This fact does not settle the later
boundary.  Exponentiation can turn a logarithmic singularity of \(G\) into a
zero of \(e^{-G}\), so we do not claim a nonremovable singularity, a natural
boundary, or nonrationality of \(\Zsol\) at \(|z|=R\).  Nor may one promote
the natural boundary of a factor in \eqref{5.1} to the infinite product:
cancellation and the domain of the product must first be controlled.

The local--global contrast is therefore precise.  Chronology changes the
analytic type of an orbit-resolved primitive base sector at the lattice scale
\(1/8\), but the leading full-system scale is the transparent branching--fibre
pole \(1/(2\cdot8)=1/16\), and the dyadic correction is strictly subleading.
